\chapter{Proof of the AR(1) representation}
\label{app:ar1-proof}

Given our model of the $\text{CO}_2$ dynamics with noise in \ref{eq:sde},
we now want to show that the process can be described as a discrete-time AR(1) process. \\
Firstly, we recognize the diffusion process as a Ornstein-Uhlenbeck, and the solution
of this process is given by \textcite[Sec.~1.13.1]{iacus2008} as 
\begin{subequations}\label{eq:ou-solution}
\begin{align}
\tilde{C}(t) &= \frac{P}{\alpha}
      + \left( \tilde{C}_0 - \frac{P}{\alpha} \right) e^{-\alpha t}
      + \sigma \int_0^t e^{-\alpha (t-u)} \, \mathrm{d}W_u
      \label{eq:ou-param1} \\
\intertext{or, in the second parametrization,}
\tilde{C}(t) &= \mu + (\tilde{C}_0 - \mu) e^{-\alpha t}
      + \sigma \int_0^t e^{-\alpha (t-u)} \, \mathrm{d}W_u ,
      \qquad \mu = \frac{P}{\alpha},
      \label{eq:ou-param2}
\end{align}
\end{subequations}
where $\tilde{C}_0$ is the initial value of the process at time $t = 0$. $\mu$ should be interpretated
as the equilibrium state of the process, which is the state that the process tends to oscillate around.


\bigskip

If we then discretize the process with a time step $\Delta t$, we can write the process as
\begin{subequations}\label{eq:ar1-proof}
\begin{align}
\tilde{C}_{t+\Delta t} &= \mu + (\tilde{C}_0 - \mu) e^{-\alpha (t + \Delta t)}
      + \sigma \int_0^{t + \Delta t} e^{-\alpha (t + \Delta t - u)} \, \mathrm{d}W_u ,
      \label{eq:ar1-full} \\
\intertext{where we can split the integral into two parts as}
\sigma \int_0^{t + \Delta t} e^{-\alpha (t + \Delta t - u)} \, \mathrm{d}W_u
  &= \sigma \biggl( e^{-\alpha \Delta t} \int_0^{t} e^{-\alpha (t-u)} \, \mathrm{d}W_u \notag \\
  &\qquad + \int_t^{t+\Delta t} e^{-\alpha (t+\Delta t-u)} \, \mathrm{d}W_u \biggr) .
      \label{eq:ar1-split} \\
\intertext{We can then write out the process as}
\tilde{C}_{t+\Delta t} &= \mu(1 - e^{-\alpha \Delta t}) + \mu e^{-\alpha \Delta t}
      + (\tilde{C}_0 - \mu) e^{-\alpha t} e^{-\alpha \Delta t} \notag \\
  \quad + \sigma e^{-\alpha \Delta t} \int_0^{t} &e^{-\alpha (t-u)} \, \mathrm{d}W_u
      + \sigma \int_t^{t+\Delta t} e^{-\alpha (t+\Delta t-u)} \, \mathrm{d}W_u
      \label{eq:ar1-expanded}
\intertext{We see then that}
\tilde{C}_{t+\Delta t} &= \biggl( \mu + (\tilde{C}_0 - \mu) e^{-\alpha t}
      + \sigma \int_0^{t} e^{-\alpha (t-u)} \, \mathrm{d}W_u \biggr) e^{-\alpha \Delta t} \notag \\
  &\quad + \mu (1 - e^{-\alpha \Delta t})
      + \sigma \int_t^{t+\Delta t} e^{-\alpha (t+\Delta t-u)} \, \mathrm{d}W_u \notag \\
    &= \tilde{C}_t e^{-\alpha \Delta t} + \mu (1 - e^{-\alpha \Delta t})
      + \sigma \int_t^{t+\Delta t} e^{-\alpha (t+\Delta t-u)} \, \mathrm{d}W_u
      \label{eq:ar1-recursion}
\end{align}
\end{subequations}
We have now shown that diffusion process is dependent on the previous emission $\tilde{C}_t$.\\ 
But for it to be an AR(1) process, we nedd to show that the drift term 
has mean 0, constant variance and is unoccorelated over time. [\textcolor{red}{MISSING REFERENCE}].\\
Let $\Delta t = c \in \mathbb{N}$ be a constant time step and given that $\alpha, t \in \mathbb{R}_{\geq 0}$, then from 
\ref{eq:ar1-recursion}, we can see that the drift term is given by
\begin{equation}
    \sigma \int_{t}^{t+\Delta t} e^{-\alpha (t+\Delta t-u)} \, \mathrm{d}W_u = 
\end{equation}
We recognize $e^{-\alpha (t+\Delta t-u)} $ as our stochastic term, which is determinitics, meaning that
that its ito intergratable as $\int_{0}^{T}e^{-2\alpha (t+\Delta t-u)} \, \mathrm{d}W_u < + \infty$. \\
Therefore, from \textcite[Thm.~1.9.1]{iacus2008},
We now know that the mean of the drit term is
\begin{equation}
    \mathbb{E}\left[\sigma \int_{t}^{t+\Delta t} e^{-\alpha (t+\Delta t-u)} \, \mathrm{d}W_u\right] = 0
\end{equation}
And that the variance is
\begin{equation}
    \text{Var}\left[\sigma \int_{t}^{t+\Delta t} e^{-\alpha (t+\Delta t-u)} \, \mathrm{d}W_u\right] = 
    \sigma^2 \int_{t}^{t+\Delta t} e^{-2\alpha (t+\Delta t-u)} \, \mathrm{d}u
\end{equation}
%Show that t to t + \delta T is 0 
%Show that the variance is independtent of t given Delta t
\begin{subequations}
\begin{align}
    \text{Var}\left[\sigma \int_{t}^{t+\Delta t} e^{-\alpha (t+\Delta t-u)} \, \mathrm{d}W_u\right]
    &= \sigma^2 \int_{t}^{t+\Delta t} e^{-2\alpha (t+\Delta t-u)} \, \mathrm{d}u \\
    &= \sigma^2 e^{-2\alpha (t+\Delta t)} \int_{t}^{t+\Delta t} e^{2\alpha u} \, \mathrm{d}u \\
    &= \sigma^2 e^{-2\alpha (t+\Delta t)} \left[\frac{1}{2\alpha} e^{2\alpha u}\right]_{t}^{t+\Delta t} \\
    &= \frac{\sigma^2}{2\alpha} e^{-2\alpha (t+\Delta t)} \left(e^{2\alpha (t+\Delta t)} - e^{2\alpha t}\right) \\
    &= \frac{\sigma^2}{2\alpha}\left(1 - e^{-2\alpha \Delta t}\right)
\end{align}
\end{subequations}
We see that the variance is independent of $t$ and 
that its constant over time. \\
We now only need to show that the covariance of the drift is 0. \\
We define that time steps as $t_1, t_2 \in \mathbb{R}_{\geq 0}$ such that $t_1 < t_2$ and that the time steps are disjoint, meaning that $[t_1, t_1 + \Delta t] \cap [t_2, t_2 + \Delta t] = \emptyset$ 
and that $t_1 + \Delta t = t_2$.
We then define $e_k$ as the drift term at time $t_k$ such that
\begin{equation}
\begin{aligned}
    \mathbb{E}\left[e_k\right] &= 0 \\ 
    \text{Var}(e_k) &= \sigma \int_{t_{k-1}}^{t_k} e^{-\alpha (t+\Delta t-u)} \, \mathrm{d}W_u = \frac{\sigma^2}{2\alpha}\left(1 - e^{-2\alpha \Delta t}\right)
\end{aligned}
\end{equation}

We now then see that the covariance of the drift term is given by, 
given that $j \neq k$
\begin{equation}
\begin{aligned}
      \text{Var}(e_k + e_j) &= \text{Var}(e_k) +  \text{Var}(e_j) + 2\text{Cov}(e_k, e_j) \\
      \Leftrightarrow 
\end{aligned}
\end{equation}



To compute $\text{Var}(e_k + e_j)$, we write both drift terms as integrals over
the same interval $[0, T]$, where $T \geq \max(t_k, t_j) + \Delta t$, by
introducing the deterministic functions
\begin{equation}
    f_k(u) = \mathbb{1}_{[t_k,\, t_k + \Delta t)}(u)\, e^{-\alpha (t_k + \Delta t - u)},
\end{equation}
such that
\begin{equation}
    e_k = \sigma \int_0^T f_k(u) \, \mathrm{d}W_u .
\end{equation}
Since $j \neq k$, the intervals $[t_k, t_k + \Delta t)$ and $[t_j, t_j + \Delta t)$
are disjoint, and therefore
\begin{equation}
    f_k(u) f_j(u) = 0 \quad \forall u \in [0, T].
\end{equation}
The sum $e_k + e_j = \sigma \int_0^T \left(f_k(u) + f_j(u)\right) \mathrm{d}W_u$
is again an Itô integral with a deterministic integrand, so by
\textcite[Thm.~1.9.1]{iacus2008}
\begin{equation}
\begin{aligned}
    \text{Var}(e_k + e_j)
    &= \sigma^2 \int_0^T \left(f_k(u) + f_j(u)\right)^2 \mathrm{d}u \\
    &= \sigma^2 \int_0^T f_k(u)^2 \, \mathrm{d}u
     + \sigma^2 \int_0^T f_j(u)^2 \, \mathrm{d}u
     + 2\sigma^2 \int_0^T f_k(u) f_j(u) \, \mathrm{d}u \\
    &= \text{Var}(e_k) + \text{Var}(e_j) + 0 .
\end{aligned}
\end{equation}
Comparing with
$\text{Var}(e_k + e_j) = \text{Var}(e_k) + \text{Var}(e_j) + 2\,\text{Cov}(e_k, e_j)$
gives
\begin{equation}
    \text{Cov}(e_k, e_j) = 0, \qquad j \neq k .
\end{equation}
Hence the drift terms have mean zero, constant variance and are uncorrelated,
and the discretized process is an AR(1) process.
